\section{反応エネルギーから電池と熱電現象へ}
\label{theory:electrochem:overview}
化学反応の係数を、熱・電子・電荷・質量の収支へつなぐ単元である。最初に反応式と物質の相を書き、次に状態量と移動量を区別する。電池の電圧と反応熱を同じ量として扱わない。

\subsection{エンタルピー差とヘスの法則}
\label{theory:electrochem:hess}
エンタルピーは状態量なので、同じ始状態と終状態を結ぶ経路の差は0である。反応をいくつかの式の和として書けるなら、反応エンタルピーも同じ係数で足せる。反応式を逆にする場合は符号を反転する。単なる結合名でなく、液体・気体・固体と温度・圧力を指定した状態同士を結ぶ必要がある\cite{openstax_hess}。
\[
 \Delta_r H^\circ=\sum_i\nu_i\Delta_f H_i^\circ .
\]
ここで生成物の係数を正、反応物を負とする。単体の標準状態の生成エンタルピーを0に選ぶのは参照の約束であり、単体がエネルギーを持たないという主張ではない。

\subsection{熱化学式の右辺とエンタルピーの符号}
\label{theory:electrochem:heat_sign}
熱化学式の「反応物＝生成物＋Q」は系が外へ放出する熱を右辺に書く慣習である。体積仕事だけの定圧過程なら系のエンタルピー変化は$-Q$になる。右辺が$-49$ kJなら吸熱で、対応する生成エンタルピーは$+49$ kJである。「生成熱」という名称だけではこの符号の約束は決まらない。候補と反応方向を合わせて確認する。

\subsection{ベンゼンの完全燃焼を組み立てる}
\label{theory:electrochem:combustion}
原典指定の液体ベンゼンについて、炭素6個と水素6個の保存から
\[
 \mathrm{C_6H_6(l)}+\frac{15}{2}\mathrm{O_2(g)}
 \longrightarrow6\mathrm{CO_2(g)}+3\mathrm{H_2O(l)}
\]
である。与えられた数値から$\Delta_rH=-3271$ kJ/molを得る。これは過去問の条件からの計算であり、外部の最新物性表へ数値を置き換えない。0.5 molなら放出熱量は1635.5 kJとなる。水を気体へ変えれば蒸発に要する分だけ熱収支が変わるため、相の指定を省かない。

\subsection{ボルン・ハーバーサイクルは符号付きの経路である}
\label{theory:electrochem:born_haber}
Na(s)と$\tfrac12$Cl$_2$(g)からNaCl(s)を作る経路を、昇華・解離・イオン化・電子獲得・結晶生成へ分ける。NaCl 1 molにはCl原子が1 mol必要なので、塩素分子の解離熱には$1/2$を掛ける。昇華、解離、イオン化は吸熱で正、Clへの電子付加と気体イオンからの格子形成は発熱で負となる。
\[
 \Delta_fH=\Delta_{\mathrm{sub}}H+\frac12D(\mathrm{Cl_2})
 +I(\mathrm{Na})-A(\mathrm{Cl})-L_{\mathrm{diss}} .
\]
ここでは$A$を電子付加の放出エネルギーの正値、$L_{\mathrm{diss}}$を格子解離の正値として定義した。電子親和力を負のエンタルピー、格子エネルギーを形成の負値で表す文献では、式の記号への代入も変わる。方向を明示せず名称だけで符号を暗記しない\cite{openstax_born_haber}。

\subsection{酸化数は電子の受け渡しを数える帳簿}
\label{theory:electrochem:oxidation_number}
単体の酸化数は0、単原子イオンは電荷に等しい。多原子イオンの各原子の酸化数の和はイオンの電荷となる。通常の酸化物ではOを$-2$と置くが、過酸化物などの例外は別に扱う。MnO$_4^-$ではMnが$+7$であり、Mn$^{2+}$への変化は酸化数の低下、すなわち還元である。酸化剤は相手を酸化して自身は還元される。実際の共有結合内の部分電荷と、形式的な酸化数は異なる\cite{mit_redox2014}。

\subsection{半反応では原子と電荷をともに保存する}
\label{theory:electrochem:half_reaction}
酸性水溶液でMnO$_4^-$をMn$^{2+}$へ還元する半反応は
\[
 \mathrm{MnO_4^-}+8\mathrm{H^+}+5e^-
 \longrightarrow\mathrm{Mn^{2+}}+4\mathrm{H_2O}
\]
である。酸素を水、水素をH$^+$、電荷を電子で釣り合わせる。酸化と還元の半反応を加えるときは、放出と受容の電子数を一致させ、電子を全反応から消す。単に「H$_2$S→S」と書かれた候補は対象元素の変化を示す略記で、原子・電荷が保存された全反応式とは区別する。

\subsection{アノードとカソードは反応で定義する}
\label{theory:electrochem:electrodes}
アノードは酸化、カソードは還元の起こる電極である。放電する電池ではアノードが負極、外から電力を与える電気分解ではアノードが正極となる。名称を常に正極又は負極と対応させない。白金は反応物の電子授受を支える不活性電極として用いられるが、すべての電池の活物質が白金という意味ではない\cite{mit_redox2014}。

\subsection{ファラデーの法則：電子molから電荷へ}
\label{theory:electrochem:faraday}
一定電流$I$を時間$t$流せば$Q=It$であり、電子の物質量は$n_e=Q/F$である。$F$の単位はC/molで、ここでいうmolは電子のmolである。物質1 molの反応に$z$ molの電子が必要なら、生成又は消費した物質の量は$Q/(zF)$となる\cite{mit_redox2014}。
\[
 m=\eta\frac{MIt}{zF} .
\]
$\eta$は対象反応に使われた電荷の割合である。副反応があると$\eta<1$となり得る。原典で全電流が電気分解に使われると明示された問題は$\eta=1$で解く。

\subsection{銅の析出では電子数2を落とさない}
\label{theory:electrochem:copper}
Cu$^{2+}+2e^-\to$Cuなので、銅0.0200 molには電子0.0400 molが要る。原典の$F=96500$ C/molでは必要電荷3860 Cである。5.36 Aで約720 s、約12分となる。ファラデー定数の現行値を持ち込まず原典の丸め値を用い、秒から分への換算を最後に行う。

\subsection{鉛蓄電池の放電と電子の係数}
\label{theory:electrochem:lead}
放電反応の全体は
\[
 \mathrm{Pb+PbO_2+2H_2SO_4\longrightarrow2PbSO_4+2H_2O}
\]
である\cite{openstax_batteries}。負極でPb(0)がPb(II)へ酸化され2電子を放出し、正極でPb(IV)がPb(II)へ還元され2電子を受け取る。反応1回の外部移動電子は2であり、両側の電子を足して4としない。硫酸は2 mol消費、水は2 mol生成するため、電子1 molあたり硫酸1 mol、水1 molが対応する。

\subsection{質量百分率の分母は使用後の液相である}
\label{theory:electrochem:solution_mass}
30質量\%の希硫酸1000 gでは、硫酸300 gと水700 gから出発する。電子1 molの放電に対応する硫酸98 gを除き、水18 gを加えると、硫酸202 g、水718 g、電解液全体920 gとなる。したがって
\[
 w=\frac{202}{202+718}\times100\%=21.96\% .
\]
減った硫酸の成分の一部は固体電極上のPbSO$_4$へ移る。閉じた電池全体の質量保存と、液相だけの質量減少は矛盾しない。蒸発、副反応、補水、充電を含む系では、それぞれの収支を追加する必要がある。

\subsection{リチウムイオンの挿入と金属の析出を分ける}
\label{theory:electrochem:lithium}
金属リチウムの負の標準還元電位と小さいモル質量は、高電圧・高い質量当たり容量につながる。一方、充電時の金属析出が針状に成長すると短絡の原因になり得る。炭素系ホストへイオンを取り込む負極は、この金属リチウムを使う経路と区別する。実用化の初期の石油コークスを、後に多用される黒鉛と歴史上同一視しない。LiCoO$_2$は正極のホスト材料の一つである\cite{nobel_lithium2019}。電極材料の内部でも電子状態は変わり、イオンが物理的に移動するだけで酸化還元がないという説明はしない。

\subsection{反応自由エネルギーと可逆起電力}
\label{theory:electrochem:gibbs}
一定温度・圧力で反応から可逆に取り出せる非体積仕事の最大値は$-\Delta_rG$である。反応を1 mol進めると$nF$の電荷を電位差$E$を通して運ぶため
\[
 \Delta_rG=-nFE .
\]
となる\cite{openstax_cell_gibbs}。ここで$n$は全反応を指定したときの電子の化学量論数である。反応式を2倍にすると$\Delta_rG$と$n$が2倍になり、起電力は変わらない。負荷時の端子電圧には内部抵抗や過電圧が含まれ、可逆起電力との差が生じる。

\subsection{ネルンスト式と二種類の「平衡」}
\label{theory:electrochem:nernst}
$\Delta_rG=\Delta_rG^\circ+RT\ln Q_r$を起電力の式へ代入すると
\[
 E=E^\circ-\frac{RT}{nF}\ln Q_r .
\]
全電池反応が化学平衡なら$\Delta_rG=0$で$E=0$となる。各電極と電解液の界面が局所平衡に達していても、二つの半電池が異なる電位を持つ開回路状態は存在する。これを全反応が完了した化学平衡と混同しない。標準電位は双方を同じ還元電位の規約で読み、カソードの値からアノードの値を引く。

\subsection{ゼーベック効果：開回路の温度積分}
\label{theory:electrochem:seebeck}
異種材料A、Bで熱電対を作ると、符号の規約を定めて起電力を
\[
 E_{AB}=\int_{T_1}^{T_2}\{S_A(T)-S_B(T)\}\,dT
\]
と表せる\cite{mit_thermoelectric2024}。$S$は熱電能、又はゼーベック係数で、一般に温度と材料へ依存する。同じ材料・同じ端点温度・開回路なら式に長さや断面積は入らない。ただし形状は抵抗、熱伝導、温度分布に影響し、負荷時の端子電圧まで形状に無関係とはならない。

\subsection{ペルティエ効果：接合部の可逆熱}
\label{theory:electrochem:peltier}
異種材料の接合部では、流れる電流に比例する可逆的な吸熱又は発熱が起こる。向きを決めて$\dot Q=(\Pi_B-\Pi_A)I$と表す。電流を反転すると吸放熱が反転するため冷却に利用できる。接合部で生じる熱と、導体全体の抵抗による発熱を区別する。

\subsection{トムソン効果：単一材料内の温度勾配}
\label{theory:electrochem:thomson}
温度勾配を持つ材料へ電流を流すと、単位体積の熱収支に電流密度と温度勾配の内積に比例する可逆項が現れる。符号規約をそろえれば$-\tau\boldsymbol j\cdot\nabla T$と書け、$\tau=T\,dS/dT$がトムソン係数である。電流又は温度勾配を反転すると符号が変わる。電流の2乗に比例する熱という説明は、この項には当てはまらない\cite{mit_thermoelectric2024}。

\subsection{ジュール熱は散逸である}
\label{theory:electrochem:joule}
抵抗$R$に電流$I$が流れると、発熱率は$I^2R$で非負である。体積密度では抵抗率$\rho_{\mathrm{el}}$に対して$\rho_{\mathrm{el}}|\boldsymbol j|^2$となる。電流方向を逆にしても熱の符号は変わらない。この反転試験により、接合部のペルティエ熱、材料内のトムソン熱、散逸によるジュール熱を区別できる。

\subsection{カー効果との名称の区別}
\label{theory:electrochem:kerr}
カー効果は電気光学的な複屈折などに関わる光学現象であり、熱電効果の名称ではない。MITの光学講義でも複屈折と電気光学変調の項目に配置される\cite{mit_kerr_calendar2005}。人名を覚える際は、原因となる量と観測する量を対応させる。熱電効果なら温度差・電流・熱流、光学効果なら電磁場と光の応答を確認する。

\section{この単元の解法の順序}
反応式・相・符号の規約を先に確定し、電子数を介して電荷と物質量を変換する。最後に濃度の分母又は指定物質量を反映する。標準電位や熱起電力の問題では、測定条件と全反応平衡の意味を添えて判断する。
